\documentclass[11pt,a4paper]{article}

\usepackage[margin=2.4cm,top=2.6cm,bottom=2.6cm]{geometry}
\usepackage{amsmath,amssymb}
\usepackage{booktabs}
\usepackage{graphicx}
\usepackage[table]{xcolor}
\usepackage{listings}
\usepackage{fancyhdr}
\usepackage{titlesec}
\usepackage{enumitem}
\usepackage[T1]{fontenc}
\usepackage[expansion=false]{microtype}
\usepackage{parskip}
\usepackage{caption}
\usepackage[hidelinks,colorlinks=true,linkcolor=accent,urlcolor=accent]{hyperref}

\definecolor{accent}{HTML}{1F4E79}
\definecolor{rule}{HTML}{B8C4CE}
\definecolor{codebg}{HTML}{F6F7F9}
\definecolor{good}{HTML}{1B6B3C}
\definecolor{warn}{HTML}{9A4B00}

\titleformat{\section}{\Large\bfseries\color{accent}}{\thesection}{0.7em}{}
\titleformat{\subsection}{\large\bfseries}{\thesubsection}{0.6em}{}
\titlespacing*{\section}{0pt}{1.6em}{0.7em}

\pagestyle{fancy}\fancyhf{}
\renewcommand{\headrulewidth}{0.4pt}
\fancyhead[L]{\small\textcolor{accent}{creditlib}}
\fancyhead[R]{\small\textcolor{accent}{Single-Name CDS Pricing Engine}}
\fancyfoot[C]{\small\thepage}

\lstset{
  basicstyle=\ttfamily\scriptsize,
  backgroundcolor=\color{codebg},
  frame=single, rulecolor=\color{rule},
  breaklines=true, showstringspaces=false,
  keywordstyle=\color{accent}\bfseries,
  commentstyle=\color{gray}\itshape,
  language=Python, xleftmargin=0pt, framexleftmargin=3pt,
  aboveskip=0.9em, belowskip=0.9em,
}

\newcommand{\keypoint}[1]{%
  \par\vspace{0.4em}\noindent
  \colorbox{codebg}{\parbox{\dimexpr\linewidth-2\fboxsep}{\small #1}}%
  \par\vspace{0.5em}}

\begin{document}

%% ------------------------------------------------------------------ title
\begin{titlepage}
\centering
\vspace*{2.4cm}
{\Huge\bfseries\color{accent} creditlib\par}
\vspace{0.7em}
{\LARGE A Single-Name CDS Pricing Engine\par}
\vspace{0.5em}
{\large Derivations, Implementation, Validation, and Application\par}
\vspace{2.4cm}
\rule{0.62\textwidth}{0.6pt}\par
\vspace{1.1em}
{\large Pearson\par}
\vspace{0.4em}
{\normalsize Prepared as a portfolio project for credit / fixed-income analyst roles\par}
\vspace{1.1em}
\rule{0.62\textwidth}{0.6pt}\par
\vspace{2.6cm}

\begin{minipage}{0.82\textwidth}\small
\textbf{Scope.} A single-name credit default swap pricing engine built from first
principles in Python. Every formula is derived before it is coded. Every component is
validated against no-arbitrage conditions and, where a comparable object exists, against
QuantLib. The library bootstraps a hazard-rate term structure from quoted CDS spreads,
prices both legs of a CDS by exact closed-form integration, and produces the standard
credit risk measures.

\medskip
\textbf{Status.} Steps 1--5 complete. 86 tests passing. Applied to a five-name
U.S.\ bank holding company cross-section.
\end{minipage}

\vfill
{\small August 2026}
\end{titlepage}

\tableofcontents
\newpage

%% ------------------------------------------------------------- 1. overview
\section{Overview and Design Philosophy}

\subsection{What the library does}

\texttt{creditlib} prices single-name credit default swaps. Given a discount curve and a
survival (credit) curve it computes the present value of both legs, the par spread, the
mark-to-market, and the standard risk measures. Run in reverse, it bootstraps the survival
curve from quoted market spreads.

The stack is Python with NumPy, SciPy, pandas and matplotlib. QuantLib is used
exclusively as an independent check --- never as a dependency of the pricing path.

\subsection{Two rules that shaped every decision}

\begin{enumerate}[leftmargin=1.4em]
\item \textbf{Derive before coding.} No formula is asserted. Every expression in this
document is obtained from a stated definition or a replication argument, so that each line
of the implementation can be defended.
\item \textbf{Validate against no-arbitrage and against QuantLib.} Structural conditions
(positivity, monotonicity, $Q(0)=1$) are enforced by construction where possible and
asserted at runtime where not. Numerical agreement with QuantLib is asserted at the
tightest tolerance the conventions permit, and where conventions differ the discrepancy is
explained and quantified rather than absorbed into a loose tolerance.
\end{enumerate}

\subsection{The design decision that made everything else work}

Both objects the pricer needs --- a discount curve and a survival curve --- have the same
mathematical form:
\[
\text{value}(t) \;=\; \exp\!\left(-\int_0^t \text{rate}(s)\,ds\right)
\]
with a \emph{piecewise-constant} instantaneous rate. For the discount curve that rate is
the instantaneous forward $f(s)$; for the credit curve it is the hazard rate $h(s)$.

Choosing piecewise-constant rates was not a convenience. It has three consequences that
each turn out to be load-bearing:

\begin{enumerate}[leftmargin=1.4em]
\item \textbf{No-arbitrage becomes a sign constraint.} A survival curve is admissible if
and only if $h \ge 0$. This is a box constraint on a scalar rather than a shape constraint
policed after the fact.
\item \textbf{The pricing integrals close in elementary form.} Because $f$ and $h$ are both
constant on the merged grid of curve knots, the protection leg and accrual integrals have
exact closed forms. No quadrature error is layered on top of the interpolation error
already accepted.
\item \textbf{The inverse problem becomes square and triangular.} With knots placed at the
quote maturities, $n$ market quotes determine exactly $n$ hazard rates, and the Jacobian is
lower triangular --- so the bootstrap is $n$ one-dimensional root-finds rather than an
$n$-dimensional nonlinear solve.
\end{enumerate}

\keypoint{Consequences (2) and (3) are why the same interpolation choice appears in both
curve classes. Had the discount curve used a spline on zero rates, the protection-leg
integral would have required numerical quadrature and the bootstrap would have lost its
clean structure.}

\subsection{Repository layout}

\begin{lstlisting}[language={}]
creditlib/
  daycount.py         ACT/360, ACT/365F, 30/360
  curves.py           PiecewiseFlatCurve -> CreditCurve
  discount_curve.py   DiscountCurve (log-linear on discount factors)
  schedule.py         IMM roll dates, ACT/360 accruals, seasoned stubs
  cds.py              legs, price_cds, par_spread, CS01 / IR DV01 / Rec01 / JTD
  bootstrap.py        survival curve from quoted spreads
  validation.py       runtime no-arbitrage assertions
tests/                79 tests
notebooks/            bank_credit_analysis.ipynb
\end{lstlisting}

\newpage
%% --------------------------------------------------- 2. survival / hazard
\section{Survival Curves and Hazard Rates}

\subsection{Setting up}

Let $\tau \ge 0$ be the default time of the reference entity, a random variable under a
risk-neutral measure. The \textbf{survival function} is
\[
Q(t) \;=\; \mathbb{P}(\tau > t), \qquad Q(0) = 1 .
\]
Assume $Q$ is absolutely continuous with density $f(t) = -Q'(t)$.

\subsection{Deriving the hazard rate}

The hazard rate is the \emph{conditional} instantaneous default intensity --- the rate of
default over the next instant, given survival to now:
\[
h(t) \;=\; \lim_{\Delta \downarrow 0}
\frac{1}{\Delta}\,\mathbb{P}\!\left(t < \tau \le t+\Delta \,\middle|\, \tau > t\right).
\]
Expanding the conditional probability, and noting
$\{t < \tau \le t+\Delta\} \subset \{\tau > t\}$ so the intersection is the numerator event:
\[
h(t) \;=\; \lim_{\Delta \downarrow 0} \frac{1}{\Delta}
\cdot \frac{Q(t) - Q(t+\Delta)}{Q(t)}
\;=\; \frac{-Q'(t)}{Q(t)}
\;=\; -\frac{d}{dt}\ln Q(t).
\]
The hazard rate \emph{is} the negative log-derivative of survival. Integrating from $0$ to
$t$ with $Q(0) = 1$:
\begin{equation}
\boxed{\;Q(t) \;=\; \exp\!\left(-\int_0^t h(s)\,ds\right) \;=\; e^{-H(t)}\;}
\end{equation}
where $H(t) = \int_0^t h(s)\,ds$ is the \textbf{cumulative hazard}.

\subsection{Two properties worth stating out loud}

\textbf{$h$ is an intensity, not a probability.} It carries units of $1/\text{time}$ and may
exceed 1. A name with $h = 3.0$ is not mis-specified; it has an expected time to default of
four months.

\textbf{Survival probabilities multiply.} $Q(t_2)/Q(t_1)$ is the probability of surviving
from $t_1$ to $t_2$ \emph{given} survival to $t_1$. This conditional structure is what makes
the bootstrap sequential: each new instrument extends the curve without disturbing what came
before.

\subsection{Piecewise-constant hazards}

Place knots at $0 = t_0 < t_1 < \dots < t_n$ and set $h(s) = h_i$ for
$s \in (t_{i-1}, t_i]$. Then
\[
H(t_j) = \sum_{i=1}^{j} h_i\,(t_i - t_{i-1}),
\qquad
H(t) = H(t_{j-1}) + h_j\,(t - t_{j-1}) \ \ \text{for } t \in (t_{j-1}, t_j].
\]
$H$ is continuous and piecewise linear, so $\ln Q$ is piecewise linear. Therefore:

\keypoint{\textbf{Piecewise-constant forward hazards and log-linear interpolation of
survival probabilities are the same assumption}, not two competing ones. This is the ISDA
CDS Standard Model convention and matches QuantLib's \texttt{BackwardFlat} interpolation.}

\subsection{Why parameterise in hazard space}

One could bootstrap $Q$ directly. With $Q = e^{-H}$ the admissibility conditions become:

\begin{center}
\begin{tabular}{ll}
\toprule
\textbf{Condition} & \textbf{Enforced by} \\
\midrule
$Q(0) = 1$ & $H(0)=0$, structural \\
$0 < Q(t) \le 1$ & $H \ge 0$, structural \\
$Q$ non-increasing & $h_i \ge 0$ --- one box constraint \\
\bottomrule
\end{tabular}
\end{center}

A rising survival curve implies $\mathbb{P}(\text{default in }(t_1,t_2]) < 0$, a genuine
arbitrage. In hazard space that is a sign constraint on a scalar; in survival space it is a
shape constraint requiring post-hoc policing. Accordingly \texttt{CreditCurve.\_\_init\_\_}
rejects negative rates, and \texttt{from\_survival\_probabilities} on a rising curve fails at
\emph{construction}, not at pricing.

\subsection{Validation}

\begin{itemize}[leftmargin=1.4em,itemsep=0.15em]
\item Flat hazard reproduces $e^{-ht}$ to $10^{-15}$.
\item $H(t) = -\ln Q(t)$ to machine zero.
\item Segment slopes of $H$ recover the input hazards to $10^{-12}$ relative.
\item Survival composes multiplicatively across sub-periods.
\item Negative hazard, and a \emph{rising} survival input, both rejected at construction.
\item \textbf{QuantLib:} survival probabilities match \texttt{ql.HazardRateCurve} to
$10^{-12}$ relative at the knots \emph{and} at interior points. Matching between knots is
the meaningful part --- it proves the interval convention $(t_{i-1}, t_i]$ agrees, not merely
the arithmetic.
\end{itemize}

\newpage
%% ----------------------------------------------------- 3. discount curve
\section{The Discount Curve}

\subsection{The discount factor is the primitive}

$D(t)$ is the price today of a contract paying 1 unit with certainty at time $t$. Three
properties, and only three:
\[
D(0) = 1, \qquad D(t) > 0, \qquad D \text{ continuous.}
\]
$D(0)=1$ is definitional. $D(t) > 0$ is the only genuine arbitrage condition --- a
non-positive price for a strictly positive payoff is free money.

\keypoint{\textbf{Nothing requires $D$ to decrease.} $D(t) > 1$ means paying 1.02 today to
receive 1.00 later, which is exactly a negative rate. EUR, JPY and CHF curves traded there
for years. The zero lower bound is an empirical regularity arising from cash storage and
custody costs, not a no-arbitrage bound. Contrast the hazard rate, where negativity
genuinely \emph{is} an arbitrage. This asymmetry is reflected in the code: the credit curve
enforces $h \ge 0$; the discount curve does not enforce $f \ge 0$.}

\subsection{Zero rates are $D$ in rate units}

A rate is meaningless without a compounding convention. Same $D(t)$, four legitimate answers:

\begin{center}
\begin{tabular}{lll}
\toprule
\textbf{Convention} & \textbf{Definition} & \textbf{Inverted} \\
\midrule
Continuous   & $D = e^{-z_c t}$              & $z_c = -\ln D / t$ \\
Annual       & $D = (1+z_a)^{-t}$            & $z_a = D^{-1/t} - 1$ \\
$m$-periodic & $D = (1 + z_m/m)^{-mt}$       & $z_m = m\left(D^{-1/(mt)} - 1\right)$ \\
Simple       & $D = (1 + z_s t)^{-1}$        & $z_s = (1/D - 1)/t$ \\
\bottomrule
\end{tabular}
\end{center}

Linked by $1 + z_a = e^{z_c}$, with $z_m \to z_c$ as $m \to \infty$. The statement
``the 5-year zero is 4.2\%'' is incomplete until the compounding \emph{and} the day count are
also stated. On the curve used later in this document, $D(5) = 0.82201223$ quotes as
3.9200\% continuous, 3.9393\% quarterly, 3.9587\% semi-annual, 3.9978\% annual, or 4.3305\%
simple --- a 411bp spread across quotations of one identical number. The library stores
$\ln D$ internally and converts on the way out, so there is a single representation and no
convention drift.

\subsection{Forward rates, derived by replication}

You wish to lend from $t_1$ to $t_2$ at a rate agreed today. Two ways to reach a known
amount at $t_2$:

\textbf{Strategy A.} Buy $1/D(t_2)$ units of the $t_2$ zero. Cost today: 1. Payoff at $t_2$:
$1/D(t_2)$.

\textbf{Strategy B.} Buy $1/D(t_1)$ units of the $t_1$ zero. Cost today: 1. At $t_1$ receive
$1/D(t_1)$ and roll to $t_2$ at a rate $f(t_1,t_2)$ agreed \emph{today}. Payoff at $t_2$:
$\frac{1}{D(t_1)}e^{f(t_1,t_2)(t_2-t_1)}$.

Both cost 1 today and pay an amount known today with certainty. Two riskless portfolios with
identical cost must have identical payoff, else one shorts the rich and buys the cheap for a
certain profit at zero net investment. Equating:
\[
\frac{1}{D(t_2)} = \frac{1}{D(t_1)}\,e^{f(t_1,t_2)(t_2-t_1)}
\qquad\Longrightarrow\qquad
\boxed{\;f(t_1,t_2) = \frac{\ln D(t_1) - \ln D(t_2)}{t_2 - t_1}\;}
\]
Under simple compounding --- what an FRA or a SOFR-style coupon actually pays --- the same
argument gives
\[
f_{\text{simple}}(t_1,t_2) = \frac{1}{\tau}\left(\frac{D(t_1)}{D(t_2)} - 1\right)
\]
where $\tau$ is the \emph{contract's} accrual fraction under \emph{its} day count, which is
generally not $t_2 - t_1$ measured in curve time.

\subsection{Instantaneous forwards: zeros are averages}

Taking $t_2 \downarrow t_1$:
\[
f(t) = -\frac{d}{dt}\ln D(t)
\qquad\Longleftrightarrow\qquad
D(t) = \exp\!\left(-\int_0^t f(s)\,ds\right),
\]
the same structure as the hazard rate. Combining with $z_c(t) = -\ln D(t)/t$:
\begin{equation}
z_c(t) = \frac{1}{t}\int_0^t f(s)\,ds
\end{equation}

\keypoint{\textbf{Zero rates are running averages of instantaneous forwards.} Forwards are
the marginal quantity, zeros the cumulative one --- so forwards are always more volatile, a
humped forward curve produces a monotone zero curve, and an interpolation that looks smooth
on zeros can be violent on forwards.}

\subsection{Why log-linear on discount factors}

Interpolate linearly in $\ln D$. For $t \in [t_i, t_{i+1}]$ with
$\theta = (t - t_i)/(t_{i+1}-t_i)$:
\[
\ln D(t) = (1-\theta)\ln D(t_i) + \theta \ln D(t_{i+1})
\qquad\Longleftrightarrow\qquad
D(t) = D(t_i)\left(\frac{D(t_{i+1})}{D(t_i)}\right)^{\theta}
\]

Four arguments, in ascending order of practical weight:

\begin{enumerate}[leftmargin=1.4em,itemsep=0.25em]
\item \textbf{Positivity is structural.} $D = \exp(\cdot) > 0$; the one genuine arbitrage
condition cannot be violated by construction.
\item \textbf{It implies piecewise-constant instantaneous forwards.} Differentiating,
$f(t)$ is the constant slope $\frac{\ln D(t_i) - \ln D(t_{i+1})}{t_{i+1}-t_i}$ on each
bucket. Log-linear DF interpolation is therefore the \emph{minimal-information} choice: the
flattest forward curve consistent with the quotes, injecting no structure the market did not
supply.
\item \textbf{It is internally consistent.} Any forward computed between two points inside
one bucket returns that bucket's forward, regardless of sampling.
\item \textbf{It localises risk.} Bump the 5Y quote: only the two adjacent buckets' forwards
move, so bucketed DV01 is clean. Under a cubic spline on zeros the perturbation propagates
across the whole curve and hedge ratios smear into untouched tenors. This is the reason
desks use it.
\end{enumerate}

\subsubsection*{What the alternatives cost}

\begin{itemize}[leftmargin=1.4em,itemsep=0.15em]
\item \textbf{Linear on $D$:} $f(t) = -D'/D = c/(a+bt)$, a hyperbola inside each bucket.
\item \textbf{Linear on zero rates:} $f(t) = z(t) + t\,z'(t)$, giving a sawtooth forward
curve with jumps at every knot. It can manufacture \emph{negative forwards from strictly
positive zero rates} whenever $z' < -z/t$, i.e.\ on any steep inversion.
\item \textbf{Cubic spline on zeros:} smooth zeros, ringing forwards, globally non-local.
\end{itemize}

Measured on the curve used in this document, comparing instantaneous forwards:

\begin{center}
\begin{tabular}{rrr}
\toprule
$t$ & log-linear DF & linear zeros \\
\midrule
3.25 & 3.8150\% & 3.8675\% \\
3.75 & 3.8150\% & 3.8325\% \\
4.25 & 3.8150\% & 3.7975\% \\
4.75 & 3.8150\% & 3.7625\% \\
5.25 & 4.1300\% & 4.0850\% \\
6.75 & 4.1300\% & 4.1750\% \\
\bottomrule
\end{tabular}
\end{center}

Linear zeros slope \emph{down} through 3--5Y and \emph{up} through 5--7Y. No quote said that.

\subsubsection*{The accepted cost}

The forward curve under log-linear DFs is a step function, discontinuous at knots. This is
irrelevant for discounting and linear products. If smooth forwards were required
(convexity-sensitive exotics), monotone-convex or tension splines would be the move, at the
price of losing property (2) and weakening (4).

\subsection{Day counts, and where money actually leaks}

Two separate jobs, routinely conflated:

\begin{enumerate}[leftmargin=1.4em]
\item \textbf{Curve time:} mapping dates to the $t$ in $D(t)$. One convention throughout.
This library uses ACT/365F, the ISDA CDS model convention.
\item \textbf{Accrual fractions:} computing $\tau$ for a specific contract's coupon,
governed by that contract's terms.
\end{enumerate}

\begin{center}
\begin{tabular}{lll}
\toprule
\textbf{Convention} & \textbf{$\tau$} & \textbf{Where it lives} \\
\midrule
ACT/360  & actual days $/$ 360 & USD \& EUR money market, SOFR, \textbf{CDS premium} \\
30/360   & $\frac{360\Delta y + 30\Delta m + \Delta d}{360}$ & US corporates, USD fixed swap leg \\
ACT/365F & actual days $/$ 365 & GBP money market, ISDA CDS curve time \\
ACT/ACT  & actual $/$ actual   & US Treasuries \\
\bottomrule
\end{tabular}
\end{center}

\textbf{The 360-day year.} ACT/360 divides by 360 but counts real days, so a calendar year
accrues $365/360 = 1.0139$ ``years'' and a rate quoted ACT/360 understates the annualised
rate: $r_{365} = r_{360}\times\frac{365}{360}$. A 5.00\% ACT/360 rate is 5.069\% ACT/365 ---
6.9bp of pure convention.

\textbf{Concrete, on 10mm over one quarter, 15 Jun $\to$ 15 Sep (92 actual days), 5.00\%:}
\[
\text{ACT/360: } \$127{,}777.78
\qquad\qquad
\text{30/360: } \$125{,}000.00
\]

\keypoint{\textbf{\$2,777.78 difference. One leg, one quarter, same rate, same dates.}
Annualised across a swap book, convention errors are a live P\&L break rather than a
rounding issue.}

\subsection{Validation}

Every compounding convention inverts its own definition to $10^{-14}$; $1+z_a = e^{z_c}$
holds; $z_m \to z_c$ monotonically ($m = 4, 52, 365, 10^4$; error $< 10^{-5}$); the
replication identity $D(t_1)e^{-f(t_1,t_2)(t_2-t_1)} = D(t_2)$ holds to $10^{-14}$, so the
derivation is tested rather than assumed; $z_c(t) = \frac1t\int_0^t f$ verified by quadrature
to $10^{-6}$; forwards constant inside a bucket to $10^{-13}$.

\textbf{QuantLib:} \texttt{ql.DiscountCurve} \emph{is} \texttt{InterpolatedDiscountCurve<LogLinear>},
so parity holds everywhere rather than only at knots --- matched at 11 points from 0.07Y to
30Y at $10^{-12}$ relative, plus \texttt{zeroRate} under both \texttt{Continuous} and
\texttt{Compounded/Annual}, plus \texttt{forwardRate} at three intervals.

\newpage
%% -------------------------------------------------------- 4. CDS pricing
\section{CDS Valuation}

\subsection{Setup and the one assumption that matters}

Value at $t=0$, notional $N$. The protection buyer pays a running coupon $s$ on dates
$t_1 < \dots < t_M$ with accrual fractions $\Delta_i$ (ACT/360). The seller pays $(1-R)N$ at
default time $\tau$ if $\tau \le T$.

\keypoint{\textbf{Default and interest rates are assumed independent.} This lets
$\mathbb{E}[D(\tau)\mathbf{1}_{\tau \le T}]$ factor into $\int D(t)(-dQ(t))$ with $D$
deterministic. It is market standard and it is \emph{wrong in exactly the states one cares
about} --- in a credit crisis defaults cluster and rates rally, so $D(\tau)$ is high
precisely when $\tau$ is small. The correction is second-order for investment-grade single
names and material for distressed. The assumption is defensible; pretending it is costless
is not.}

\subsection{Protection leg}

The payoff is $(1-R)$ at the random time $\tau$, contingent on $\tau \le T$:
\[
\text{PV}_{\text{prot}}
= (1-R)\,\mathbb{E}\!\left[D(\tau)\mathbf{1}_{\tau\le T}\right]
= (1-R)\int_0^T D(t)\,\big(-dQ(t)\big).
\]
Here $-dQ(t)$ is the unconditional probability of defaulting in $[t, t+dt)$. Converting to
hazard form using $Q = e^{-H}$, so $Q'(t) = -h(t)Q(t)$:
\begin{equation}
\boxed{\;\text{PV}_{\text{prot}} = (1-R)\int_0^T D(t)\,h(t)\,Q(t)\,dt\;}
\end{equation}
Read directly: $h(t)\,dt$ is the chance of defaulting in the next instant given survival,
$Q(t)$ the chance of getting there, $D(t)$ discounts it.

\subsubsection*{Evaluating it exactly}

Most implementations Riemann-sum this integral. They need not. On any interval $[a,b]$ where
$f$ and $h$ are \emph{both} constant --- guaranteed on the merged grid of discount-curve and
credit-curve knots --- write $D(t)Q(t) = D(a)Q(a)e^{-(f+h)(t-a)}$ and set $c = f+h$:
\[
\int_a^b D h Q\,dt
= h\,D(a)Q(a)\!\int_0^{b-a}\!\!e^{-cu}du
= \frac{h}{c}\,D(a)Q(a)\left[1 - e^{-c(b-a)}\right],
\]
and since $D(b)Q(b) = D(a)Q(a)e^{-c(b-a)}$ this collapses to
\begin{equation}
\boxed{\;\int_a^b D\,h\,Q\,dt \;=\; \frac{h}{f+h}\Big[D(a)Q(a) - D(b)Q(b)\Big]\;}
\end{equation}

\keypoint{Summing over sub-intervals gives an \textbf{exact} protection leg. The only error
is the curve interpolation already chosen upstream --- no quadrature error is layered on top.
This is precisely why Steps 1 and 2 both used piecewise-flat rates.}

Sanity check: setting $f = 0$ makes the factor 1 and the sum telescopes to
$(1-R)[Q(0)-Q(T)] = (1-R)(1-Q(T))$, the undiscounted expected loss.

\subsection{Premium leg}

\subsubsection*{Scheduled coupons}
\[
\text{PV}_{\text{coupon}} = s\sum_{i=1}^{M}\Delta_i\,D(t_i)\,Q(t_i)
\]

\subsubsection*{Accrual on default --- why the term exists}

Premium accrues \emph{continuously} but pays \emph{quarterly}. If default lands at
$\tau \in (t_{i-1}, t_i)$, the buyer has enjoyed protection from $t_{i-1}$ to $\tau$ without
paying for it, since the coupon was not yet due. The contract requires that stub at
settlement.

\keypoint{Without this term, a CDS would be worth a different amount depending on
\emph{where in the coupon cycle it is traded}. Protection bought the day after a coupon date
would carry up to 90 days of free protection in the default state; bought the day before, a
full quarter's premium for one day of cover. The standardised 100bp/500bp fixed-coupon
contract would not be fungible across trade dates. \textbf{The accrual term is what makes the
instrument standardisable} --- not a rounding refinement.}

\[
\text{PV}_{\text{accr}} = s\sum_i \frac{\Delta_i}{t_i - t_{i-1}}
\int_{t_{i-1}}^{t_i}(t - t_{i-1})\,D(t)\,h(t)\,Q(t)\,dt
\]
The prefactor converts a curve-time offset into an accrual fraction; it is exact whenever the
day count is proportional within a period, true for ACT/360 and ACT/365F up to day-rounding.

The same closed form is available. On $[a,b]$ inside a period starting at $p$, with $c=f+h$
and $L = b-a$, using $\int_0^L u e^{-cu}du = \frac{1-e^{-cL}(1+cL)}{c^2}$:
\[
\int_a^b (t-p)DhQ\,dt
= h\,D(a)Q(a)\left\{\frac{1 - e^{-cL}(1+cL)}{c^2} + (a-p)\,\frac{1-e^{-cL}}{c}\right\}
\]

\subsection{The risky annuity and mark-to-market}
\[
A \;=\; \underbrace{\sum_i \Delta_i D(t_i)Q(t_i)}_{\text{coupons}}
\;+\; \underbrace{\sum_i \frac{\Delta_i}{t_i-t_{i-1}}\int(t-t_{i-1})DhQ\,dt}_{\text{accrual}},
\qquad \text{PV}_{\text{prem}} = s\cdot A
\]
$A$ does not depend on $s$. For the protection buyer (long protection, short credit):
\[
\text{MTM}_{\text{buyer}} = \text{PV}_{\text{prot}} - s\cdot A
\]
Under the post-2009 standard contract $s$ is fixed at 100bp or 500bp and the MTM is exchanged
upfront, so this is literally the upfront payment.

\subsection{Par spread: a ratio, not a solve}
\begin{equation}
\text{MTM}=0 \;\Longrightarrow\;
\boxed{\;s_{\text{par}} = \frac{\text{PV}_{\text{prot}}}{A}\;}
\end{equation}
Because the protection leg does not involve $s$ and $A$ does not either, the par spread is a
quotient --- no root-finding. This matters precisely because the bootstrap \emph{does} require
a solve: there the unknown is the hazard rate, not the spread. It also gives
\[
\text{MTM}_{\text{buyer}} = (s_{\text{par}} - s_{\text{contract}})\cdot A ,
\]
from which CS01 follows.

\subsection{Clean versus dirty par spread}

Under the post-2009 standard the accrual period begins at the IMM date preceding the trade,
so the buyer pays a \emph{full} first coupon and is rebated the accrued portion. A par spread
computed from the raw annuity therefore includes a stub of coupon for which no protection is
provided.

\keypoint{The market convention is the \textbf{clean} par spread, which nets the accrued
rebate out of the annuity. On a 1-year contract that stub can be 16\% of the annuity.
Bootstrapping from the dirty spread overstates the front hazard by roughly 18bp --- a 20\%
error at the front of the curve. \texttt{par\_spread(clean=True)} is the default.}

\subsection{The credit triangle, exactly}

With flat $h$, flat $r$, and \emph{continuously paid} premium:
\[
A_{\text{cont}} = \int_0^T e^{-(r+h)t}dt = \frac{1-e^{-(r+h)T}}{r+h},
\qquad
\text{PV}_{\text{prot}} = (1-R)h\cdot\frac{1-e^{-(r+h)T}}{r+h}
\]
Dividing:
\begin{equation}
s_{\text{par}} = h(1-R) \qquad \textbf{exactly, independent of } r \textbf{ and } T
\end{equation}

\keypoint{The credit triangle is not an approximation in this limit --- it is an identity. The
approximation error in practice comes entirely from \emph{discrete} premium payment. This is
the primary validation target for the whole pricer.}

\subsection{Risk measures --- four different things}

\begin{center}
\begin{tabular}{p{2.3cm}p{5.6cm}p{6.4cm}}
\toprule
\textbf{Measure} & \textbf{Definition} & \textbf{Note} \\
\midrule
RPV01 ($A$) & PV of 1 unit of spread, survival-contingent & Not a risk number; a \emph{unit}. Quoted per unit notional. \\[0.3em]
CS01 & $\Delta$MTM for $+1$bp on the par spread curve, re-bootstrapped & $\approx A\cdot N\cdot 10^{-4}$ to first order. Second-order gap because moving spreads moves $Q$, which moves $A$. \\[0.3em]
IR DV01 & $\Delta$MTM for $+1$bp on the discount curve & Small: the legs largely offset. \\[0.3em]
JTD & $(1-R)N - \text{accrued} - \text{MTM}$ & Discontinuous. Not a derivative. \\
\bottomrule
\end{tabular}
\end{center}

\subsection{Hand-checked example}

Flat $r = 5\%$ continuous, flat $h = 2\%$, $R = 40\%$, $T = 1$, one annual payment, no
accrual-on-default. Every quantity is derivable on paper in three lines:
\[
\text{PV}_{\text{prot}} = 0.60\cdot\tfrac{0.02}{0.07}\cdot(1-e^{-0.07}),
\qquad A = e^{-0.07},
\qquad s_{\text{par}} = \tfrac{\text{PV}}{A}
\]

\begin{center}
\begin{tabular}{lrrr}
\toprule
& \textbf{Engine} & \textbf{Hand} & \textbf{Difference} \\
\midrule
PV$_{\text{prot}}$ & 0.011589630873 & 0.011589630873 & $1.7\times10^{-18}$ \\
RPV01              & 0.932393819906 & 0.932393819906 & 0 \\
$s_{\text{par}}$   & 124.299739 bp  & 124.299739 bp  & --- \\
\bottomrule
\end{tabular}
\end{center}

\subsection{Convergence to the credit triangle}

Target $h(1-R) = 120$bp:

\begin{center}
\begin{tabular}{rrr}
\toprule
\textbf{Frequency} & \textbf{No accrual} & \textbf{With accrual} \\
\midrule
1   & 124.2997 bp & 123.0401 bp \\
4   & 121.0562 bp & 120.7525 bp \\
52  & 120.0808 bp & 120.0577 bp \\
365 & 120.0117 bp & 120.0083 bp \\
\bottomrule
\end{tabular}
\end{center}

Monotone, and inside 0.01bp at daily frequency. Verified independent of $r$ and $T$ across
nine $(r,T)$ combinations.

\subsubsection*{What the accrual term actually does}

A sharper claim than ``it improves accuracy''. Over a period of length $\delta$ the discrete
coupon undervalues the continuous one by $\approx (f+h)\frac{\delta^2}{2}DQ$, while the
accrual integral contributes $\approx h\frac{\delta^2}{2}DQ$. It therefore cancels
\emph{exactly the hazard component} of the discretisation error and leaves the discounting
component. Prediction: the residual error ratio tends to $f/(f+h) = 0.05/0.07 = 0.7143$.
Measured at quarterly frequency: $0.7525/1.0562 = 0.7125$. This is now a test.

\subsection{The 8bp gap, decomposed}

A QuantLib comparison (flat 4\%/3\%, 5Y IMM) prices at 172.03bp, not the 180bp the credit
triangle predicts. That gap is contract economics, not error, and factors cleanly:

\begin{enumerate}[leftmargin=1.4em]
\item \textbf{ACT/360 versus ACT/365F.} Every $\Delta_i$ is $365/360$ too large relative to
curve time, inflating the annuity. Measured ratio: 0.9863014 against
$360/365 = 0.9863014$ --- nine digits.
\item \textbf{Seasoned stub.} Accrual began 20 Jun; protection runs from 17 Aug. The annuity
carries $58/360$ of a coupon buying no protection --- 3.59\% of RPV01.
\end{enumerate}
\[
180.90 \times \tfrac{360}{365} \times (1 - 0.0359) = 172.014\text{bp}
\quad\text{versus}\quad 172.027\text{bp priced}
\]
A residual of 0.013bp. Trading the same maturity \emph{on} an IMM date recovers 178.42bp:
wedge 2 vanishes, wedge 1 remains.

\subsection{Worked example: 5Y, 10mm notional, 100bp coupon}

\begin{center}
\begin{tabular}{lr}
\toprule
Protection leg PV & 390,711.35 \\
Premium leg PV & 466,061.41 \\
\quad coupons & 465,222.72 \\
\quad accrual & 838.69 \\
RPV01 (per unit) & 4.660614 \\
MTM & $-$75,350.06 \\
Par spread & 83.8326 bp \\
Accrued (rebate) & 16,111.11 \\
Upfront points & $-$0.7535\% \\
\midrule
CS01 (reprice) & 4,689.72 \\
CS01 (analytic) & 4,660.61 \ ($-$0.62\%) \\
IR DV01 & 6.54 \ (0.14\% of CS01) \\
Rec01 ($+$1pt) & $-$6,511.86 \\
JTD & 6,059,238.95 \ (1,292$\times$ CS01) \\
\bottomrule
\end{tabular}
\end{center}

Buying 100bp protection on an 84bp name loses 75k, as it should. Accrual is 0.18\% of the
premium leg in PV --- nearly nothing --- while moving the par spread meaningfully. IR DV01 is
0.14\% of CS01: the legs almost fully offset.

\keypoint{\textbf{JTD is 1,292$\times$ CS01.} A book can be flat CS01 and still be destroyed
by a single default. That ratio is why desks limit both, and why an index hedge that
neutralises spread risk leaves idiosyncratic default risk entirely untouched.}

\newpage
%% ------------------------------------------------------- 5. bootstrapping
\section{Bootstrapping the Survival Curve}

\subsection{The inverse problem is underdetermined until parameterised}

Forward: given $h(\cdot)$, compute $s_{\text{par}}(T)$. Backward: given $n$ quotes, find
$h(\cdot)$. Infinitely many hazard functions reproduce $n$ quotes.

\keypoint{\textbf{The market determines $n$ numbers, not a function.} The bootstrapped curve
is \emph{quotes plus an interpolation assumption}, never ``the market's view of default''.}

Piecewise-constant $h$ with knots placed \emph{exactly at the quote maturities} gives $n$
unknowns for $n$ equations: square system, unique solution, zero residual. Knots anywhere else
turn this into least-squares fitting, a different problem with different failure modes.

\subsection{Why sequential: the Jacobian is triangular}

Define the residual $\Phi_k = \text{PV}_{\text{prot},k} - s_k A_k$. The $k$-th CDS matures at
$T_k$, so both legs integrate against $Q(t)$ only for $t \le T_k$, and $Q$ on that range
depends only on $h_1,\dots,h_k$. Hence
\[
\frac{\partial \Phi_k}{\partial h_j} = 0 \quad \text{for all } j > k,
\qquad
J = \begin{pmatrix}
\ast & 0 & 0 & 0 & 0\\
\ast & \ast & 0 & 0 & 0\\
\ast & \ast & \ast & 0 & 0\\
\ast & \ast & \ast & \ast & 0\\
\ast & \ast & \ast & \ast & \ast
\end{pmatrix}
\]

Lower triangular, so the solution is forward substitution: an $n$-dimensional nonlinear
system collapses to $n$ one-dimensional root-finds. Four consequences, in ascending order of
practical weight:

\begin{enumerate}[leftmargin=1.4em,itemsep=0.25em]
\item \textbf{Robustness.} Bracketed 1-D root-finding on a monotone function essentially
always converges. Multivariate Newton needs a good starting point and can diverge into
negative hazards.
\item \textbf{Exact reproduction.} $n$ equations, $n$ unknowns --- the fit is exact, not
least-squares.
\item \textbf{Economic meaning.} $h_k$ is the \emph{marginal} default intensity over
$(T_{k-1},T_k]$ implied by the incremental information in the $k$-th quote. It is a forward
quantity, like a forward rate. The 5Y quote does not describe aggregate 0--5Y default risk;
it describes 3--5Y risk given what 1Y and 3Y already said.
\item \textbf{Error localisation.} A bad 10Y quote cannot corrupt $h_1$. Errors propagate
forward, never backward.
\end{enumerate}

\subsection{Why solve $\Phi$ rather than the par-spread residual}

Both have the same root and the same monotonicity, but the ratio form degenerates: at the
first tenor $A \to 0$ as $h$ grows, so $s_{\text{par}}$ is unbounded (reaching 3.01 at $h=5$
on a 1Y contract) while $\Phi$ stays bounded within $[-0.012, 0.589]$ and is measurably less
curved. Both failure modes remain detectable in the $\Phi$ form: $\Phi_k(0) > 0$ catches the
inverted-curve case and bracket-expansion failure catches the ceiling.

\subsection{Why each 1-D solve is well posed}

Raising $h_k$ raises $\text{PV}_{\text{prot}}$ (the contribution over $(T_{k-1},T_k]$ rises
monotonically toward the certain-default limit $D(T_{k-1})Q(T_{k-1})$) and lowers the coupon
part of $A$ (every $Q(t_i)$ for $t_i > T_{k-1}$ falls). Both push $\Phi$ the same way, so
$\Phi$ is strictly increasing and has at most one root.

One honest caveat: the \emph{accrual} part of $A$ also rises with $h$, so denominator
monotonicity is not unconditional. That term is $O(h\delta^2)$ and is dominated over any
realistic range. It is verified numerically over $h \in [10^{-6}, 5]$ rather than asserted as
a theorem.

\subsection{Failure modes fall out of the bracketing}

\subsubsection*{Spread above the certain-default ceiling}

As $h_k \to \infty$, default becomes certain immediately after $T_{k-1}$:
\[
s^{\max}(T_k) \;=\;
\frac{\text{PV}_{\text{prot}}(T_{k-1}) + (1-R)\,D(T_{k-1})Q(T_{k-1})}
     {A(T_{k-1}) + \varepsilon_{\text{accr}}}
\]
Finite for every $k \ge 2$ (unbounded for $k=1$, where $A \to 0$). A quote above $s^{\max}$
cannot be reproduced by any non-negative hazard: it prices protection above the maximum
possible loss given the earlier quotes. That is an arbitrage in the quote set, and the
bootstrap must say so rather than silently returning garbage.

\subsubsection*{Negative implied forward hazard on an inverted curve}

Inversion alone is normal --- distressed names invert routinely, because near-term default
risk dominates. What fails is inversion steep enough that reproducing $s_k < s_{k-1}$ requires
$h_k < 0$, which means $Q$ \emph{rises} over $(T_{k-1},T_k]$: a negative forward default
probability. The arbitrage is to sell $T_{k-1}$ protection and buy $T_k$ protection, being
paid to own a strictly non-negative claim over the window.

A rough threshold from the credit triangle: the implied forward hazard turns negative when
$s_k T_k \lesssim s_{k-1}T_{k-1}$, i.e.\ when the spread$\times$maturity product falls.

\subsubsection*{Recovery inconsistency and the upfront problem}

$R$ is an assumption, not a quote. Pushing $R$ toward 1 collapses loss given default, so
$s^{\max}$ falls and previously fine quotes become unbootstrappable. Distressed names are the
acute case: the market stops quoting running spread and quotes \emph{points upfront against a
fixed coupon}, precisely because the par-spread representation degenerates. Bootstrapping a
2000bp name from a ``par spread'' is a modelling error before it is a numerical one.

\subsection{Results}

Quotes at 1/3/5/7/10Y of 55/82/105/118/130bp, $R = 40\%$:

\begin{center}
\begin{tabular}{rrrrrrr}
\toprule
\textbf{Tenor} & \textbf{Quoted} & \textbf{Repriced} & \textbf{Error} & \textbf{Hazard} & \textbf{Fwd spread} & \textbf{Survival} \\
\midrule
1.0  & 55.0000 & 55.0000 & $-2.6\times10^{-14}$ & 92.35  & 55.41  & 0.992263 \\
3.0  & 82.0000 & 82.0000 & $-1.7\times10^{-14}$ & 158.48 & 95.09  & 0.961264 \\
5.0  & 105.0000 & 105.0000 & $+3.5\times10^{-14}$ & 239.62 & 143.77 & 0.916283 \\
7.0  & 118.0000 & 118.0000 & $-5.2\times10^{-14}$ & 263.71 & 158.22 & 0.869146 \\
10.0 & 130.0000 & 130.0000 & $+1.7\times10^{-14}$ & 282.38 & 169.43 & 0.798488 \\
\bottomrule
\end{tabular}
\end{center}

Maximum repricing error: $4.44\times10^{-16}$ relative --- two units in the last place of a
float64. Verified two independent ways: par spread equality, and pricing each quoted CDS at
its own coupon off the bootstrapped curve ($|\text{MTM}| < 10^{-7}$ on 10mm).

\keypoint{Note the hazard column against the forward-spread column. The 5Y quote is 105bp,
but the marginal intensity over 3--5Y implies a 143.77bp forward spread. \textbf{Bootstrapped
hazards are forward quantities.} Reading $h_5$ as ``the 5Y default intensity'' is the standard
error.}

\subsection{Triangularity, made testable}

\texttt{test\_locality\_of\_the\_triangular\_system} bumps the 10Y quote by 50bp and asserts
$h_1..h_4$ are unchanged using \texttt{np.array\_equal} --- bit-for-bit, not approximately.
\texttt{test\_subset\_bootstrap\_matches\_prefix} bootstraps only the first three quotes and
asserts exact equality with the first three hazards of the full run. That \emph{is} forward
substitution.

\subsection{Feasible quote windows}

\begin{center}
\begin{tabular}{rrrr}
\toprule
\textbf{Tenor} & \textbf{Floor} & \textbf{Ceiling} & \textbf{Quoted} \\
\midrule
1Y  & 0.00 bp   & 119,429.26 bp & 55.00 \\
3Y  & 16.95 bp  & 6,576.34 bp   & 82.00 \\
5Y  & 50.36 bp  & 1,985.56 bp   & 105.00 \\
7Y  & 78.17 bp  & 1,154.55 bp   & 118.00 \\
10Y & 88.80 bp  & 807.05 bp     & 130.00 \\
\bottomrule
\end{tabular}
\end{center}

The floor rises and the ceiling collapses moving out the curve: each new quote is constrained
on both sides by everything preceding it. The 1Y has no floor (nothing precedes it) and an
effectively unbounded ceiling.

\subsection{The recovery degeneracy}

Same quotes, two recovery assumptions, both repricing to $4\times10^{-16}$:
\[
\frac{h(R=40\%)}{h(R=20\%)} = [1.333,\ 1.334,\ 1.337,\ 1.340,\ 1.343]
\qquad\text{against}\qquad \frac{0.8}{0.6} = 1.333
\]
Implied 5Y cumulative default probability: 8.37\% at $R=40\%$, 6.34\% at $R=20\%$.

\keypoint{\textbf{CDS spreads identify expected loss $h(1-R)$, not $h$.} A two-point spread in
implied default probability from an assumption nobody observes. Any sentence beginning ``the
CDS market is pricing an $X\%$ chance of default'' is really ``\dots given someone's recovery
assumption''.}

\newpage
%% ---------------------------------------------------------- 6. application
\section{Application: Five U.S.\ Bank Holding Companies}

\subsection{What this analysis can and cannot claim}

\keypoint{\textbf{One quote per name is not a term structure.} Each name supplies a single 5Y
spread, so the bootstrap has one degree of freedom and returns a \emph{flat} hazard. Every
point away from 5Y is the piecewise-flat assumption, not market information. The chart marks
the 5Y point as observed and renders everything else dashed. Obtaining 1/3/5/7/10Y quotes per
name would produce a genuine term structure --- the engine already handles it; the constraint is
data, not model.}

Two further constraints carried into every conclusion below:

\begin{itemize}[leftmargin=1.4em,itemsep=0.15em]
\item \textbf{Risk-neutral, not actuarial.} These probabilities embed a default risk premium,
a liquidity premium, and correlation compensation. They run roughly 8--12$\times$ historical
default rates for comparable ratings.
\item \textbf{Indicative spreads, not executable marks.} Post-2014 single-name U.S.\ bank CDS
is materially less liquid than pre-crisis, with much flow in CDX.NA.IG.
\end{itemize}

\subsection{Why 40\% recovery, and why it is the weakest input}

40\% is the \textbf{ISDA quoting convention} for senior unsecured obligations of
developed-market corporates. It makes quotes comparable. It is \emph{not} an estimate of what
these instruments would recover. For U.S.\ bank holding companies the convention is doing more
work than usual:

\begin{enumerate}[leftmargin=1.4em,itemsep=0.25em]
\item \textbf{Holdco senior is the bail-in layer.} Single-name CDS on U.S.\ banks references
the \emph{holding company}. Under Dodd-Frank Title II Orderly Liquidation Authority and the
single-point-of-entry resolution strategy, holdco senior unsecured is explicitly TLAC --- it
exists to be written down or converted so the operating subsidiaries continue. Its recovery in
the scenario that triggers a CDS credit event is by design low, not the corporate-average 40\%.
\item \textbf{The one clean data point is 8.625\%.} Lehman Brothers senior CDS settled at
8.625 cents in the October 2008 auction --- roughly a fifth of the convention.
\item \textbf{The empirical base is thin.} Bank senior recoveries are governed by resolution
mechanics never tested at scale on a U.S.\ G-SIB.
\end{enumerate}

Practically: 40\% is defensible \emph{as a convention} and is the right choice for consistency
with how the quotes were formed. It is not defensible as a forecast.

\subsection{Results}

\begin{center}
\begin{tabular}{lrrrrr}
\toprule
\textbf{Name} & \textbf{5Y spread} & \textbf{Hazard} & \textbf{1Y PD} & \textbf{5Y PD} & \textbf{10Y PD} \\
\midrule
JPMorgan        & 58 bp & 96.2 bp  & 0.96\% & 4.70\% & 9.17\% \\
Bank of America & 63 bp & 104.5 bp & 1.04\% & 5.09\% & 9.92\% \\
U.S.\ Bancorp   & 67 bp & 111.1 bp & 1.11\% & 5.41\% & 10.52\% \\
Wells Fargo     & 72 bp & 119.4 bp & 1.19\% & 5.80\% & 11.26\% \\
Citigroup       & 85 bp & 141.0 bp & 1.40\% & 6.81\% & 13.15\% \\
\bottomrule
\end{tabular}
\end{center}

\begin{figure}[htbp]
\centering
\includegraphics[width=\textwidth]{bank_implied_pd.png}
\caption{Left: risk-neutral cumulative default probability at $R=40\%$. Solid to the quote,
dashed beyond --- the dashed portion is the flat-hazard assumption, not market data. Right:
implied 5Y PD under four recovery assumptions.}
\end{figure}

\subsection{Three structural findings}

\subsubsection*{The recovery assumption cannot touch the ranking}

With a uniform $R$ across names, $h = s/(1-R)$ is a monotone transform of the spread, so the
ordering is invariant --- recovery moves levels only. Reordering requires \emph{name-specific}
recovery, and the required differential is large:

\begin{center}
\begin{tabular}{lr}
\toprule
\textbf{Ties JPMorgan at $R=40\%$ if its own recovery is} & \\
\midrule
Bank of America & 34.83\% \\
U.S.\ Bancorp   & 30.69\% \\
Wells Fargo     & 25.52\% \\
Citigroup       & 12.07\% \\
\bottomrule
\end{tabular}
\end{center}

Citigroup only ties JPMorgan under a 28-point recovery differential. \textbf{State the ordering
with confidence; state the levels with a recovery caveat attached.}

\subsubsection*{Rates are irrelevant here}

A 600bp swing in the discount curve moves Citigroup's implied 5Y PD from 6.8300\% to 6.7807\%
--- five basis points of probability. The Step 4 IR DV01 result, confirmed end to end.

\subsubsection*{The entire cross-sectional signal is 27bp}

The three tightest names sit inside 9bp of each other; JPM$\to$BAC is 0.394pp of PD on a 5bp
spread gap. Whether JPM, BAC and USB are genuinely distinguishable on this data is a question
about bid/ask, not about credit.

\subsection{Risk-neutral versus actuarial}

\begin{center}
\begin{tabular}{lr}
\toprule
Aa-rated historical 5Y cumulative default rate & 0.31\% \\
A-rated historical 5Y cumulative default rate & 0.56\% \\
Baa-rated historical 5Y cumulative default rate & 1.55\% \\
\midrule
Market-implied range & 4.70\% -- 6.81\% \\
Ratio to A-rated actuarial & 8.4$\times$ -- 12.2$\times$ \\
\bottomrule
\end{tabular}
\end{center}

Roughly an order of magnitude above realised default experience. That wedge is the risk
premium --- compensation for jump-to-default, illiquidity, and the fact that bank defaults
arrive correlated with everything else in a portfolio. It is not a forecast error; it is why
risk-neutral PD is a \emph{pricing} quantity rather than an \emph{actuarial} one.

\subsection{Market-implied ranking against fundamental ranking}

\subsubsection*{Where they agree: Citigroup}

The market puts Citi 27bp wide of JPMorgan --- a 47\% spread premium and the only gap in the
set that clears any plausible bid/ask. A thin-PPNR-coverage thesis and the market's ordering
point the same direction.

Precision matters about what that agreement is worth: Citi's spread is also consistent with
lower through-cycle ROTCE, ongoing transformation and consent-order costs, and a larger EM
footprint. Concluding ``the market is pricing my PPNR thesis'' is over-attribution. The
defensible claim is narrower --- \emph{the market independently ranks Citi last, and PPNR
coverage is one of a small number of explanations consistent with that}.

\subsubsection*{Where they disagree: Bank of America}

A held-to-maturity tail thesis says BAC carries the largest unrealised securities loss position
in the peer set relative to tangible capital. The market prices BAC as the \emph{second-safest}
name --- 5bp inside U.S.\ Bancorp and 9bp inside Wells Fargo. On this data the market does not
price the HTM tail as a credit event driver at all.

\subsubsection*{The reframe}

The instinct is to argue the market is wrong. The stronger move is to notice that CDS and a
fundamental screen are answering \emph{different questions}, and that both can be right.

A CDS prices the probability and severity of a \emph{credit event}. HTM losses do not cause a
credit event through the income statement. They cause one only through a \textbf{funding
channel}: forced liquidation of the HTM book to meet outflows, crystallising the mark and
destroying capital. That is the SVB mechanism, and it is a liquidity-plus-concentration
failure, not a securities-portfolio failure.

For BAC that channel is weak: a very large, granular, heavily insured retail deposit base, a
G-SIB LCR and NSFR regime, discount window access, and a post-March-2023 policy precedent for
term lending against par. The market appears to be saying the HTM position is a \emph{book
value and NII problem, not a solvency problem} --- that it suppresses tangible book,
AOCI-adjusted CET1 and forward earnings power without moving the probability of resolution.

\keypoint{If that is right, the HTM thesis is not contradicted by 63bp --- it is
\textbf{mispriced in a different instrument}. A view that BAC's HTM tail impairs earnings power
and capital return is an \emph{equity} view, or a preferred/AT1 view. Senior CDS is the wrong
expression vehicle, and the fact that CDS does not reflect it is evidence about instrument
selection rather than evidence against the thesis.

\medskip
\emph{``The market and I are not disagreeing about Bank of America. We are pricing different
questions --- and my HTM thesis belongs in the capital structure below senior, where the loss
actually lands.''}}

\subsubsection*{The falsifiable version}

This framing generates a test rather than an opinion. If the HTM tail is genuine credit risk
that senior CDS is missing, it should appear first in funding-sensitive instruments: BAC's
preferred/AT1 spreads widening relative to peers, the senior--sub CDS basis widening, or
deposit beta and uninsured-deposit share deteriorating faster than the peer set. If those are
flat while the equity discount persists, the market's separation is correct and the thesis is
an earnings thesis wearing credit clothing.

\subsection{Highest-value extensions}

\begin{enumerate}[leftmargin=1.4em]
\item \textbf{Obtain a term structure.} 1/3/5/7/10Y quotes per name turn five numbers into
five curves, and the \emph{shape} is where a funding thesis would show up --- a concern that is
real but slow should steepen the front end relative to peers. A single 5Y point structurally
cannot express that.
\item \textbf{Add the sub-senior structure.} Senior/sub CDS basis and preferred spreads make
the reframe above testable rather than rhetorical.
\item \textbf{Pull current holdco ratings} as an independent third ranking. A
market/fundamental/agency three-way comparison is materially more persuasive than a two-way.
\end{enumerate}

\newpage
%% -------------------------------------------------------- 7. validation
\section{Validation Summary}

86 tests, all passing. QuantLib is used only as an independent oracle.

\begin{center}
\begin{tabular}{lrp{8.2cm}}
\toprule
\textbf{Module} & \textbf{Tests} & \textbf{Coverage} \\
\midrule
\texttt{curves.py} & 15 & Analytic identities, no-arbitrage rejection, QuantLib \texttt{HazardRateCurve} parity at $10^{-12}$ \\
\texttt{discount\_curve.py} & 17 & Convention inversions, replication identity, bucket-constant forwards, QuantLib \texttt{DiscountCurve} parity at $10^{-12}$ \\
\texttt{cds.py} + \texttt{schedule.py} & 33 & Two hand-checked closed forms, credit-triangle convergence, exact-vs-quadrature, QuantLib \texttt{MidPointCdsEngine} within 1bp \\
\texttt{bootstrap.py} & 21 & Machine-precision repricing, triangularity, failure modes, recovery degeneracy, QuantLib \texttt{PiecewiseFlatHazardRate} parity \\
\bottomrule
\end{tabular}
\end{center}

\subsection{Selected assertions}

\begin{itemize}[leftmargin=1.4em,itemsep=0.2em]
\item \textbf{Exact integration versus brute force.} The closed-form protection integral is
compared against a Riemann sum at grid sizes 20k, 200k and 2M. Because the integrand is
discontinuous at every hazard knot, trapezoid converges slowly --- so the test asserts
\emph{convergence toward} the closed form rather than agreement at one grid size. That is the
correct direction of the claim: the closed form is exact and the Riemann sum is the
approximation.
\item \textbf{Exact versus midpoint.} The two methods agree to 0.0023bp on a 5Y par spread and
are asserted \emph{not} identical --- if they were, one implementation would be wrong.
\item \textbf{Locality.} \texttt{np.array\_equal} on hazards after a 50bp bump to the 10Y
quote. Bit-for-bit, not approximate.
\end{itemize}

\subsection{Defects found and corrected}

Documented because the process is part of the artifact.

\begin{enumerate}[leftmargin=1.4em,itemsep=0.4em]
\item \textbf{Inverted accrual direction.} An assertion that ACT/360 produces a higher quoted
rate than ACT/365F failed. ACT/360 yields the \emph{larger} year fraction
($91/360 = 0.25278$ against $0.24932$), so for identical growth the quoted simple rate is
\emph{lower}. ``Divide by 360'' feels as though it should shrink the fraction. It does not.

\item \textbf{RPV01 scaled by notional.} The risky annuity was being reported multiplied by
notional rather than per unit. Mathematically harmless (it cancels in the par spread) but
against market convention and confusing on output. Corrected to the per-unit figure.

\item \textbf{Wrong maturity convention.} Schedules measured tenor from the \emph{trade date},
making every contract up to a full quarter too long --- a ``5Y'' maturing 20 Sep 2031 against
QuantLib's 20 Jun 2031. The CDS2015 convention measures tenor from the \emph{accrual start}
(the preceding IMM date). This produced a 66bp error in the bootstrapped 1Y hazard.

\item \textbf{Bootstrapping off the dirty par spread.} On a seasoned schedule the annuity
carries a coupon stub buying no protection. On a 0.84Y contract that stub is 16\% of the
annuity, inflating $h_1$ by 18bp --- 20\% wrong at the front of the curve. Netting the accrued
rebate collapsed the QuantLib gap:

\begin{center}
\begin{tabular}{lrrrr}
\toprule
& \textbf{Dirty} & \textbf{Clean} & \textbf{QuantLib} & \textbf{Clean $-$ QL} \\
\midrule
1Y  & 110.38 & 92.44  & 92.14  & $+0.31$ bp \\
3Y  & 162.46 & 158.44 & 158.47 & $-0.03$ bp \\
10Y & 284.22 & 282.31 & 282.28 & $+0.04$ bp \\
\bottomrule
\end{tabular}
\end{center}

\item \textbf{Closure mutation in the bootstrap loop.} The residual function closed over the
hazard list, which was appended to before the diagnostic re-evaluation --- lengthening the
array out from under the closure. Caught by 19 simultaneous test failures.

\item \textbf{Clean par spread against dirty mark-to-market.} Introducing the clean par
spread (defect 4) netted the accrued rebate out of the annuity for $s_{\text{par}}$ but not
for the MTM, which continued to use the dirty annuity. On a seasoned schedule the identity
$\text{MTM} = (s_{\text{par}} - s)\,A\,N$ therefore failed by exactly the accrued rebate, and
pricing \emph{at} the par spread returned $-\text{accrued}$ rather than zero. Invisible to the
suite because every MTM test used the synthetic \texttt{flat\_schedule}, where
\texttt{accrued\_fraction()} is zero and the two conventions coincide. Surfaced by a user
running the engine on a seasoned IMM schedule, not by inspection. Corrected by making
\texttt{mtm} the clean market value, exposing \texttt{stream\_pv} for the raw cashflow view
(which jump-to-default correctly nets against, since on default the stub is owed rather than
rebated), and adding seven regression tests on a deliberately seasoned fixture.

\item \textbf{Incorrect claim about recovery sensitivity.} A draft chart asserted that the
recovery assumption drove the \emph{ranking}. It does not: with uniform $R$, $h = s/(1-R)$ is
a monotone transform of the spread and the ordering is invariant. Recovery drives the
\emph{level}. Caught by running the sensitivity rather than reasoning about it.
\end{enumerate}

\keypoint{Defects 3, 4 and 6 were each surfaced by a validation step, not by inspection.
Defect 4 was diagnosed in a single step because the Step 4 day-count decomposition had already
isolated the accrued-rebate wedge as a named quantity.}

\newpage
%% ------------------------------------------------------- 8. limitations
\section{Limitations and Roadmap}

\subsection{Known simplifications}

\begin{enumerate}[leftmargin=1.4em,itemsep=0.3em]
\item \textbf{No business-day adjustment.} Payment dates should roll Following against a real
holiday calendar. This moves a handful of cashflows by 1--3 days and is the residual in the
QuantLib comparisons ($\sim$0.3bp at 1Y, under 0.07bp beyond). Mechanical to add; omitted as a
scoping choice and stated in the module docstring so the omission is auditable.

\item \textbf{Rates and default assumed independent.} Standard, and wrong in stress. Material
for distressed names.

\item \textbf{Deterministic recovery.} $R$ is a fixed input, not a distribution, and no
recovery/default correlation is modelled.

\item \textbf{Running-spread quotation only.} The post-2009 standard contract quotes points
upfront against a fixed 100bp or 500bp coupon. For distressed names the par-spread
representation degenerates as the annuity collapses. An upfront bootstrap is the natural next
addition.

\item \textbf{No curve-shape alternatives.} Piecewise-flat hazards are the only
parameterisation. The bootstrapped curve is quotes \emph{plus} that assumption.
\end{enumerate}

\subsection{Roadmap}

\begin{enumerate}[leftmargin=1.4em,itemsep=0.2em]
\item Points-upfront quotation and conversion to and from par spread.
\item Bucketed CS01 computed off a full re-bootstrap, exploiting the locality property.
\item Business-day adjustment against a real calendar.
\item Senior/subordinated basis, for the capital-structure analysis in Section 6.
\item Index (CDX) decomposition against single-name constituents.
\end{enumerate}

\subsection{Reproducing this document}

\begin{lstlisting}[language=bash]
pip install numpy scipy pandas matplotlib pytest QuantLib
PYTHONPATH=. python -m pytest tests/ -q          # 79 passed
jupyter notebook notebooks/bank_credit_analysis.ipynb
\end{lstlisting}

\vspace{2em}
\hrule
\vspace{0.8em}
{\small\itshape
All figures in this document are produced by the code in the repository. Bank CDS spreads are
indicative and supplied as inputs; the discount curve is illustrative. Historical default rates
should be verified against the current Moody's Annual Default Study before external citation.
Nothing here is investment advice.}

\end{document}
